\chapter[1972 --- Edelman et al.]{1972: Gerald M. Edelman, Rodney R. Porter}
\label{ch:1972_Edelman}

\section*{Main Contribution}

\textit{Determined the chemical structure of antibodies, explaining how these molecules recognize and neutralize specific pathogens with extraordinary precision.}

\section{Official Citation}

for their discoveries concerning the chemical structure of antibodies

\section{Background and Historical Context}

The 1972 Nobel Prize in Physiology or Medicine was awarded to Gerald M. Edelman and Rodney R. Porter for their discoveries concerning the chemical structure of antibodies. This prize recognized two related breakthroughs that collectively elucidated the molecular structure of antibodies---the proteins produced by the immune system that recognize and neutralize foreign pathogens.

\subsection{The Antibody Puzzle}

By the mid-20th century, it was well established that antibodies play a central role in the immune response. Antibodies are proteins produced by B lymphocytes (a type of white blood cell) that can recognize and bind to specific foreign molecules called antigens. However, the structure of antibodies---how their atoms are arranged in three-dimensional space---was still unknown.

Understanding the structure of antibodies was essential for understanding how they recognize and bind to antigens, and for developing new immunotherapies. The work of Edelman and Porter would provide the first detailed understanding of antibody structure.

\subsection{Gerald M. Edelman}

Gerald M. Edelman was born on July 1, 1929, in New York City. He studied at the University of Pennsylvania, where he earned his bachelor's degree in 1950, and at the Rockefeller Institute for Medical Research, where he earned his Ph.D. in 1960. After completing his doctorate, Edelman joined the Rockefeller Institute, where he would spend most of his career.

Edelman was a skilled biochemist and immunologist who was interested in the structure of antibodies. His work on the amino acid sequence of antibodies would earn him the Nobel Prize.

\subsection{Rodney R. Porter}

Rodney R. Porter was born on October 8, 1917, in Liverpool, England. He studied at the University of Liverpool, where he earned his bachelor's degree in 1939, and at the University of Cambridge, where he earned his Ph.D. in 1948. After completing his doctorate, Porter worked at the National Institute for Medical Research in London and later at St. Mary's Hospital Medical School, where he would spend most of his career.

Porter was a skilled biochemist who was interested in the structure and function of antibodies. His work on the fragmentation of antibodies would earn him the Nobel Prize.

\subsection{The Structure of Immunoglobulins}

Antibodies, also known as immunoglobulins, are Y-shaped proteins that consist of four polypeptide chains: two identical heavy chains and two identical light chains. The chains are linked together by disulfide bonds, and the antigen-binding site is located at the tips of the Y.

Understanding this structure was essential for understanding how antibodies recognize and bind to antigens. The work of Edelman and Porter would provide the first detailed understanding of antibody structure.

\section{The Discovery/Contribution}

The work of Edelman and Porter on antibody structure involved the determination of the complete amino acid sequence of antibodies and the elucidation of their three-dimensional structure.

\subsection{Porter's Work on Antibody Fragmentation}

Rodney Porter's primary contribution was the fragmentation of antibodies using enzymes, which revealed the basic structure of the molecule.

\subsubsection{Papain Digestion}

Porter discovered that the enzyme papain could cleave antibodies into three fragments:

\begin{itemize}
\item Two identical Fab fragments (Fragment antigen-binding): Each fragment contains one antigen-binding site
\item One Fc fragment (Fragment crystallizable): This fragment does not bind to antigens but mediates other effector functions of antibodies
\end{itemize}

\subsubsection{Significance of the Discovery}

Porter's work on antibody fragmentation was significant because it revealed the basic structure of antibodies. The Fab fragments contain the antigen-binding sites, while the Fc fragment mediates effector functions such as complement activation and binding to immune cells.

\subsection{Edelman's Work on Antibody Sequence}

Gerald Edelman's primary contribution was the determination of the complete amino acid sequence of an antibody molecule.

\subsubsection{The Experimental Approach}

Edelman used chemical methods to reduce the disulfide bonds that hold the antibody chains together and then separated the heavy and light chains. He then determined the amino acid sequence of each chain using traditional protein sequencing methods.

\subsubsection{Key Results}

Edelman discovered that:

\begin{itemize}
\item The light chains consist of approximately 220 amino acids, with a variable region (V region) at the N-terminus and a constant region (C region) at the C-terminus
\item The heavy chains consist of approximately 450 amino acids, with a variable region at the N-terminus and a constant region at the C-terminus
\item The variable regions of both the heavy and light chains contribute to the antigen-binding site
\item The constant regions determine the effector functions of the antibody
\end{itemize}

\subsubsection{Significance of the Discovery}

Edelman's work on antibody sequence was significant because it revealed the genetic basis of antibody diversity. The variable regions of antibodies are encoded by different gene segments that can be rearranged in different combinations, generating an enormous diversity of antigen-binding sites.

\subsection{The Three-Dimensional Structure}

Together, the work of Edelman and Porter revealed the three-dimensional structure of antibodies:

\begin{itemize}
\item The antibody molecule is Y-shaped, with two identical Fab arms and one Fc stem
\item The antigen-binding sites are located at the tips of the Fab arms
\item The Fc stem mediates effector functions
\item The molecule is flexible, allowing the Fab arms to move relative to each other
\end{itemize}

\section{Impact on Modern Medicine}

The work of Edelman and Porter on antibody structure has had a significant impact on modern medicine. Their discoveries have contributed to the understanding and treatment of infectious diseases, autoimmune diseases, and cancer, and have led to the development of numerous immunotherapies.

\subsection{Understanding Infectious Diseases}

The understanding of antibody structure has been essential for understanding how the immune system fights infectious diseases. By knowing how antibodies recognize and bind to pathogens, researchers can design vaccines that stimulate the production of effective antibodies.

\subsection{Monoclonal Antibodies}

The understanding of antibody structure has also contributed to the development of monoclonal antibodies---laboratory-produced molecules that can mimic the immune system's ability to fight harmful pathogens. Monoclonal antibodies are now used to treat a wide range of diseases, including cancer, autoimmune diseases, and infectious diseases.

\subsection{Cancer Immunotherapy}

The understanding of antibody structure has also contributed to the development of cancer immunotherapy. Monoclonal antibodies can be designed to target specific molecules on the surface of cancer cells, allowing the immune system to recognize and destroy the cancer cells. This approach has been successful in treating a variety of cancers, including breast cancer, lung cancer, and lymphoma.

\subsection{Autoimmune Diseases}

The understanding of antibody structure has also contributed to the treatment of autoimmune diseases. In autoimmune diseases, the immune system produces antibodies that attack the body's own tissues. Understanding the structure of these autoantibodies has led to the development of new treatments that can suppress the immune response without causing excessive immunosuppression.

\subsection{Transplant Medicine}

The understanding of antibody structure has also contributed to transplant medicine. Antibodies play a major role in transplant rejection, and understanding their structure has led to the development of new strategies for preventing transplant rejection.

\subsection{Diagnostic Tests}

The understanding of antibody structure has also contributed to the development of diagnostic tests. Antibody-based diagnostic tests are now widely used to detect and diagnose a wide range of diseases, including infectious diseases, autoimmune diseases, and cancer.

\section{Legacy}

The legacy of Gerald Edelman and Rodney Porter extends far beyond their Nobel Prize-winning work on antibody structure. Their discoveries fundamentally changed our understanding of immunology and have contributed to numerous advances in medicine and biotechnology.

\subsection{Foundation for Immunology}

The work of Edelman and Porter established the foundations of modern immunology. Their discoveries about antibody structure provided the first detailed understanding of how antibodies recognize and bind to antigens. The principles they established continue to guide research in immunology today.

\subsection{Advances in Medicine}

The understanding of antibody structure has contributed to numerous advances in medicine, including the development of monoclonal antibodies, cancer immunotherapy, and new treatments for autoimmune diseases. These advances have improved the lives of millions of patients and continue to drive progress in these fields.

\subsection{Biotechnology Applications}

The understanding of antibody structure has also contributed to the biotechnology industry. Monoclonal antibodies are now one of the fastest-growing segments of the pharmaceutical industry, with billions of dollars in annual sales. The development of new antibody-based therapies continues to be a major focus of pharmaceutical research.

\subsection{Antibody Engineering}

The understanding of antibody structure has also led to the development of antibody engineering---the design of antibodies with improved properties for therapeutic use. These include:

\begin{itemize}
\item Humanized antibodies: Antibodies that have been modified to reduce their immunogenicity in humans
\item Bispecific antibodies: Antibodies that can bind to two different antigens simultaneously
\item Antibody-drug conjugates: Antibodies linked to toxic drugs that can specifically target and kill cancer cells
\item Single-chain variable fragments (scFvs): Small antibody fragments that can be produced in bacteria and are easier to manufacture than full-length antibodies
\end{itemize}

\subsection{Inspiration for Future Researchers}

The work of Edelman and Porter continues to inspire researchers in immunology and biochemistry. Their success in using elegant experiments to answer fundamental questions about antibody structure demonstrates the power of basic research to improve human health. Their example has motivated generations of scientists to pursue careers in immunology and biochemistry.

\subsection{Recognition and Honors}

In addition to the Nobel Prize, Edelman and Porter received numerous other honors and awards for their work. Edelman received the National Medal of Science in 1989 and was elected to the National Academy of Sciences. Porter received the Order of the British Empire in 1972 and was elected to the Royal Society.

Their contributions to immunology and medicine are remembered and celebrated by researchers and clinicians around the world. Their discoveries continue to influence our understanding of antibody structure and have led to numerous advances in medicine and biotechnology, ensuring that their legacy will endure for generations to come.


\section{Modern Developments}

Edelman and Porter's antibody structure work has led to modern antibody engineering. Monoclonal antibodies (Köhler and Milstein, Nobel 1984) have become major therapeutics. Antibody-drug conjugates (ADCs) deliver chemotherapy specifically to cancer cells. Bispecific antibodies engage immune cells to attack tumors. Nanobodies (single-domain antibodies from camels) offer new therapeutic possibilities. Understanding of antibody structure has enabled rational design of therapeutic antibodies.
